In this setting, SAE-based feature decomposition works effectively on a purely visual Vision Transformer: SAEs trained on ViT-B/16 MLP activations recover sparse, interpretable features across four depths, and CLIP-based external evaluation, with an abstention margin, provides a viable labelling strategy without internal text supervision. Across depth we find a coherent hierarchy: early features (Layer~2) are uncertain texture detectors, a mid layer (Layer~9) is net-inhibitory, and late layers (Layers~10--11) concentrate positive causal weight on cleanly monosemantic object features, the strongest being \textit{American lobster} (Feature~1321, $+0.56\%$ relative logit drop, 95\% bootstrap CI $[+0.17,+1.04]$). By estimating every causal effect over a held-out set with bootstrap confidence intervals, we avoid the single-image extrapolation that limited earlier ViT analyses. A controlled comparison with a self-supervised DINOv2-B backbone is equally informative: despite a $99.5\%$ linear-probe accuracy, the identical SAE recipe collapses to dead latents at all but one depth, showing that sparse-dictionary recipes calibrated on supervised activations do not transfer to self-supervised ones without re-tuning. Limitations include the finite, ImageNet-biased CLIP vocabulary, the restriction to MLP outputs, the absence of explicit patch-to-\texttt{[CLS]} aggregation tracing, and the visual rather than large-panel human assessment of monosemanticity. Natural extensions include per-backbone SAE re-tuning for a fair DINOv2 monosemanticity study, and attention-level circuit analysis (e.g., activation patching on attention heads) to trace classification-specific computation into the \texttt{[CLS]} token.
